%% The MIT License (MIT)
%%
%% Copyright (c) 2015 Daniil Belyakov

\documentclass[]{mcdowellcv}

\usepackage{amsmath}
\usepackage{hyperref}
\geometry{top=0.3in}

\name{Aditya Kulkarni}
\address{\hspace{0.2\textwidth}\href{https://www.kulkarniaditya.com}{kulkarniaditya.com} \linebreak \href{https://www.linkedin.com/in/adityakulkarni244/}{linkedin.com/in/adityakulkarni244}}
\contacts{(281)-876-7066 \linebreak adityakulkarnius@gmail.com \linebreak \href{https://github.com/Aditya-k24}{github.com/Aditya-k24}}

\begin{document}

\makeheader

\begin{cvsection}{Employment}
  \begin{cvsubsection}{Product Engineer Intern}{Isomer AI}{May 2025 -- May 2026}
    \begin{itemize}
      \item Shipped 10+ production AI features: multi-model LLM pipelines (Claude, GPT-4o, Gemini), MCP tool orchestration, Temporal.io durable workflows, and Auth0 RBAC; 923 commits across 7 production deployments.
      \item Own backend stack (PostgreSQL, Redis, Docker, Terraform, Sentry APM); maintain CI/CD and resolve production incidents.
    \end{itemize}
    \vspace{-1.5pt}
  \end{cvsubsection}

  \begin{cvsubsection}{Graduate Research Assistant}{NC State University}{April 2025 -- Present}
    \begin{itemize}
      \item Deploy genomic platform on GCP (Cloud Run) with JBrowse 2 genome browser; support 100+ researchers, 500GB+ datasets.
      \item Improve ML model accuracy by 12\% via hyperparameter tuning and time-series analysis on biological sensor data.
    \end{itemize}
    \vspace{-1.5pt}
  \end{cvsubsection}

  \begin{cvsubsection}{Software Developer Intern}{Mphasis (HP Company)}{Jun 2024 -- Aug 2024}
    \begin{itemize}
      \item Built 5+ AI workflows including health claim automation (OCR, NLP, ML); improved document processing accuracy by 25\%.
      \item Automated AWS and GCP deployments via CI/CD; integrated Grafana monitoring and led model red-teaming for reliability.
    \end{itemize}
    \vspace{-1.5pt}
  \end{cvsubsection}

  \begin{cvsubsection}{AI/ML Intern}{Capgemini}{Jun 2023 -- Aug 2023}
    \begin{itemize}
      \item Led SDLC for data-driven ticket analysis system using NLP classification; reduced average resolution time by 15\%.
      \item Built Kafka + Spark + Hadoop pipelines for batch and real-time data ingestion; cut ticket-processing latency by 20\%.
    \end{itemize}
    \vspace{-1.5pt}
  \end{cvsubsection}
\end{cvsection}

\begin{cvsection}{Education}
  \begin{cvsubsection}{Raleigh, NC}{North Carolina State University}{Aug 2024 -- May 2026}
    \begin{itemize}
      \item Master of Science in Computer Science
      \item Courses: Artificial Intelligence, Operating Systems, Algorithms, Data Science, Object Oriented Development, Software Engineering.
    \end{itemize}
    \vspace{-1.5pt}
  \end{cvsubsection}

  \begin{cvsubsection}{Mumbai, India}{Mumbai University}{Aug 2020 -- May 2024}
    \begin{itemize}
      \item Bachelor of Technology in Computer Science; CGPA: 9.32/10. Courses: Data Warehousing, Big Data Infra, DS \& Algorithms, Machine Learning, Blockchain.
    \end{itemize}
    \vspace{-1.5pt}
  \end{cvsubsection}
\end{cvsection}

\begin{cvsection}{Projects}
  \begin{cvsubsection}{LLM Eval Harness}{Vendor-Neutral LLM Evaluation Framework}{2025}
    \begin{itemize}
      \item Built vendor-neutral LLM evaluation harness in Python comparing Claude, GPT-4o, and Gemini across 2,000 frozen labeled examples spanning factual recall, multi-hop reasoning, and FEVER claim verification tasks.
      \item Tracked ExactMatch, TokenF1, hallucination rate, abstention accuracy, p50/p95 latency, and cost-per-call; applied paired bootstrap CIs, McNemar tests, and Holm correction across 3 repeated runs per model.
    \end{itemize}
    \vspace{-1.5pt}
  \end{cvsubsection}

  \begin{cvsubsection}{CortexQ}{Kubernetes LLM Inference Observability Platform}{2026}
    \begin{itemize}
      \item Built scale-to-zero Kubernetes LLM inference platform: KEDA autoscaler on Redis queue-depth metrics, Kopf CRD operator for multi-provider routing across Claude, GPT-4o, and Gemini, with circuit breaker and EWMA load balancing.
      \item Instrumented OpenTelemetry distributed tracing, Prometheus, and Grafana to expose p95 latency under 500-VU k6 load; Helm + ArgoCD GitOps; 88 pytest tests in CI.
    \end{itemize}
    \vspace{-1.5pt}
  \end{cvsubsection}

  \begin{cvsubsection}{VectorLift}{Semantic Search \& ML Evaluation Pipeline}{2026}
    \begin{itemize}
      \item Built production semantic search on MS MARCO: Elasticsearch BM25 baseline, sentence-transformers bi-encoder fine-tuned with in-batch negatives, and MiniLM cross-encoder re-ranker; benchmarked 6 pipeline configurations with NDCG@10, MRR@10, MAP, and paired bootstrap significance tests.
      \item Designed Kafka-driven indexing pipeline for event-based corpus ingestion into Elasticsearch and Qdrant; served search via FastAPI with per-stage latency tracking and Streamlit experiment dashboard.
    \end{itemize}
  \end{cvsubsection}
\end{cvsection}

\begin{cvsection}{Languages and Technologies}
  \begin{cvsubsection}{}{}{}
    \begin{itemize}
      \item \textbf{Languages:} Python, TypeScript, Java, JavaScript, SQL, Bash
      \item \textbf{AI/ML \& Observability:} LLM evaluation, GenAI pipelines, RAG, OpenTelemetry, Prometheus, Grafana, Sentry
      \item \textbf{MLOps \& Cloud:} Kubernetes, Docker, Helm, ArgoCD, KEDA, AWS, GCP, Terraform, GitHub Actions
      \item \textbf{Data \& Messaging:} Kafka, Elasticsearch, Qdrant, PostgreSQL, Redis, MongoDB, FastAPI, Node.js
    \end{itemize}
  \end{cvsubsection}
\end{cvsection}

\begin{cvsection}{Additional Experience and Awards}
  \begin{cvsubsection}{}{}{}
    \begin{itemize}
      \item \textbf{IEEE Published:} Co-authored paper on CV-Based Cybersecurity Threat Detection (Springer icSoftComp 2023); presented to 500+ attendees.
      \item \textbf{National Hackathon Winner:} Won or ranked Top 10 in 6+ national hackathons spanning ML, systems, and developer tools; built 3 complete MVPs end-to-end within 48-hour sprints, each with a working demo.
      \item \textbf{Technical Writing:} Published on LLM model collapse, AI agent security, and audio AI at \href{https://www.kulkarniaditya.com/blog}{kulkarniaditya.com/blog}
    \end{itemize}
  \end{cvsubsection}
\end{cvsection}

\end{document}
